\section{Limitations}
\label{sec:limitations}

% Scope of the claim: one card, one architecture, one corpus, one scale.
% Which code paths carry a correctness argument and which do not; which
% comparisons are confounded and what would separate them.

\paragraph{The scale is small.}
$11{,}061{,}531$ parameters on a $6$\,GB card, picked so all four arms fit at
once: a model large enough to fill the card would have priced out the low-rank
arm and the five-member ensemble. What transfers is the ratio between arms on
one implementation, not the absolute cost. The character-level vocabulary of
$27$ keeps the token embedding negligible, so no result here covers a posterior
over a large embedding table. The position embedding is sized for a
$1024$-token context while training runs at $64$, so about $3\%$ of the
parameters never see a data gradient and are held at the prior by the KL term
alone, while still costing posterior state in three of the four arms.

\paragraph{Nothing is converged.}
$100{,}000$ steps at $256$ tokens per step is $25.6$M tokens against a $90$M
token training split, $28\%$ of one epoch. Every number is what a method does
at a fixed step budget, and the ranking between arms is free to change with
more steps.

\paragraph{Step times are lower bounds, not controlled measurements.}
The runs span two sessions on a card with a cooling fault. It throttled to
$46$--$51$\,W at $87$--$89$\,\textdegree C against a board budget of
$65$--$80$\,W, and the low-rank arm's step time drifted from $52$ to $65$\,ms
within a single run with no change in the work being done. The two sessions did
not throttle alike, so the arms were not clocked under matched conditions. A
gap of a few milliseconds is not a result, and what a healthy card of this
class would deliver is not measured here.

\paragraph{$n = 1$, and the noise floor is unknown.}
One training run per arm, no error bars, no seed sweep. The runs are also not
reproducible in principle: \texttt{encoder\_backward},
\texttt{layernorm\_backward} and \texttt{attention\_backward} accumulate
gradients with \texttt{atomicAdd}, so the order of floating-point additions
varies and two identical invocations diverge. How large that divergence grows
over $100{,}000$ steps is not measured, so which differences between arms
exceed it is unknown.

\paragraph{The off-distribution stream is synthetic.}
Uniform-random token ids, not shifted text. Sample disagreement has to rise
where the input carries no structure, and an arm that fails that is not
capturing model uncertainty. It supports no comparison with the shift
literature \citep{ovadia2019trust}, for which a second text domain would be the
minimum.

\paragraph{Noise scale and posterior quality are confounded.}
Loss tracks the injected noise scale monotonically across arms, so an arm that
samples less aggressively looks better on predictive metrics for a reason
unrelated to its posterior. Separating the two needs a sweep over the noise
scale within each arm and a comparison of curves rather than of single points,
which runs against saved checkpoints.

\paragraph{What each code path's correctness argument covers.}
Each kernel was validated against the CPU reference implementation, but that
reference is a single-block model and says nothing about the block loop at
$L > 1$. That gap is closed by a finite-difference gradient check on a two-block
configuration, run per parameter group and per block, which agrees to
$\sim\!10^{-4}$ relative, the same error the single-block configuration shows.
Vadam's optimiser kernels carry two more checks: a double-precision host
reference for the update, and a distributional test of the sampler over
$20{,}000$ draws covering the marginal variances and the absence of
off-diagonal correlation. Those checks are why Section~\ref{sec:quality}
attributes that arm's posterior to $N s$ against $\lambda$ rather than to a
bug. None of this covers the fused FlashAttention kernels beyond the CPU
reference, and none of it is a proof.

\paragraph{The arms are not identical in what they cover.}
The low-rank arm places its posterior over the four GEMM weight matrices per
block and leaves embeddings, LayerNorm parameters and biases deterministic,
where the other three are Bayesian over every parameter. It covers $96\%$ of
the model rather than all of it, which flatters its cost. The ensemble
evaluates at $S = E = 5$ by construction while the Bayesian arms evaluate at
whatever $S$ is requested, so a mixture metric compared across arms at
different $S$ compares different sample counts.
